\begin{abstract}
We address the EPIC-KITCHENS-100 Multi-Instance Retrieval (MIR) challenge through three progressively stronger approaches: (i)~a JPoSE-based zero-fine-tuning baseline that exploits part-of-speech embeddings; (ii)~a score-level ensemble of three pre-trained retrieval models (JPoSE, MI-MM and MMEN) with optional graph-diffusion re-ranking on visual and textual k-nearest-neighbour graphs; and (iii)~end-to-end fine-tuning of an AVION ViT-L/14 dual encoder with the Symmetric Multi-Similarity (SMS) loss, which directly consumes the soft-label relevancy matrices, enhanced with horizontal-flip and longer-clip test-time augmentation. Across six iterations, average nDCG on the official Codabench leaderboard rises from \textbf{53.53} (JPoSE alone) to \textbf{55.82} (ensemble + re-ranking) and finally to \textbf{69.68} (AVION ViT-L + SMS, 16~epochs, with TTA), an absolute gain of \textbf{+16.15}~nDCG. We analyse where each gain originates, whether encoder capacity, soft-label supervision, test-time augmentation and additional epochs, and outline a model-soup-style ensemble of differently-trained ViT-L+SMS instances as future work, motivated by the compute and offline-evaluation constraints encountered. Our code will be released at: \href{https://github.com/ISIS-4825-Assignments/Multi-Instance-Retrieval-EK-100}{https://github.com/ISIS-4825-Assignments/Multi-Instance-Retrieval-EK-100}
\end{abstract}